% D5
% Auto-generated from markdown source
% origin: ai

\documentclass[12pt, a4paper]{article}
\usepackage[utf8]{inputenc}
\usepackage[T1]{fontenc}
\usepackage{amsmath, amssymb, amsfonts, amsthm}
\usepackage{textcomp}
\usepackage{geometry}
\usepackage{hyperref}
\usepackage{booktabs}
\usepackage{graphicx}
\usepackage{fancyhdr}
\usepackage{tcolorbox}
\usepackage{enumitem}

\geometry{margin=1in}
\pagestyle{fancy}
\fancyhf{}
\rhead{D5}
\lhead{\thepage}

\newtheorem{definition}{Definition}[section]
\newtheorem{theorem}{Theorem}[section]
\newtheorem{corollary}{Corollary}[section]
\newtheorem{proposition}{Proposition}[section]
\newtheorem{lemma}{Lemma}[section]
\newtheorem{conjecture}{Conjecture}[section]

\newcommand{\calI}{\mathcal{I}}
\newcommand{\calP}{\mathcal{P}}
\newcommand{\calE}{\mathcal{E}}
\newcommand{\calS}{\mathcal{S}}
\newcommand{\calH}{\mathcal{H}}
\newcommand{\calO}{\mathcal{O}}
\newcommand{\calF}{\mathcal{F}}
\newcommand{\calU}{\mathcal{U}}
\newcommand{\calV}{\mathcal{V}}
\newcommand{\calA}{\mathcal{A}}
\newcommand{\calG}{\mathcal{G}}
\newcommand{\calB}{\mathcal{B}}
\newcommand{\calC}{\mathcal{C}}
\newcommand{\calL}{\mathcal{L}}
\newcommand{\calM}{\mathcal{M}}
\newcommand{\calX}{\mathcal{X}}
\newcommand{\calY}{\mathcal{Y}}
\newcommand{\calW}{\mathcal{W}}
\newcommand{\calT}{\mathcal{T}}
\newcommand{\calZ}{\mathcal{Z}}
\newcommand{\calR}{\mathcal{R}}
\newcommand{\calQ}{\mathcal{Q}}
\newcommand{\calD}{\mathcal{D}}
\newcommand{\calN}{\mathcal{N}}
\newcommand{\calK}{\mathcal{K}}

\newcommand{\Tr}{\text{Tr}}
\newcommand{\E}{\mathbb{E}}
\newcommand{\diff}{\mathrm{d}}
\newcommand{\KL}{D_{\mathrm{KL}}}
\newcommand{\argmax}{\arg\max}
\newcommand{\argmin}{\arg\min}
\newcommand{\hbarF}{\hbar}

\newtcolorbox{thmbox}[1]{
  colback=blue!5!white, colframe=blue!75!black, fonttitle=\bfseries,
  title={#1}, breakable
}
\newtcolorbox{defbox}[1]{
  colback=green!5!white, colframe=green!75!black, fonttitle=\bfseries,
  title={#1}, breakable
}

\begin{document}

\title{D5}
\author{RALPH Loop}
\date{}
\maketitle

\begin{abstract}
This paper presents D5, formalizing the underlying principles using information-theoretic methods.

\textbf{Keywords}: D5, information theory, formalization.
\end{abstract}

\tableofcontents
\section{D5}




\#\#



\#\#

\textbf{-} \(\rho_\theta\) Bures

\[d_B(\rho_1, \rho_2) = \sqrt{2(1 - \text{Tr}(\rho_1^{1/2}\rho_2\rho_1^{1/2}))}\]

Fisher
\[g_{ij} = \text{Re}\left[\text{Tr}(\rho_\theta L_i L_j)\right]\]

\(L_i\) SLD

\textbf{*}*** \(S(\rho) = -\text{Tr}(\rho \log \rho)\)  Betti

\[\beta_1(\mathcal{P}_\theta) \leq C \cdot S(\rho_\theta) + \log(d)\]

\(d\)


\#\#

\subsubsection{ 1 Bures-}


\textbf{-}
\[\lim_{\hbar \to 0} d_B(\rho_\theta, \rho_{\theta+\delta\theta}) = \|\theta - (\theta + \delta\theta)\|_{g_{FI}}\]

\textbf{ Fisher Braunstein-Caves, 1994}
\[g_{ij}^{Q} = \lim_{\hbar \to 0} g_{ij}^{classical} + \mathcal{O}(\hbar)\]

**** D5-D1-1/D1-2


\subsubsection{ 2-Betti}


\textbf{ Betti } n  Schmidt  \(|\psi\rangle = \sum_i \lambda_i |i\rangle_A |i\rangle_B\)

\[\beta_1(\mathcal{P}_\theta) \leq \mathcal{O}\left(n \cdot e^{S_E/2}\right)\]

\(S_E = -\sum_i \lambda_i^2 \log \lambda_i^2\) ****

****
\begin{itemize}
\item ""
\item D1-2 \(\|T\|_{HS}/\|T\|_{op}\)  \(S(\rho)\)
\end{itemize}



\subsubsection{ 3}


\textbf{Tishby et al.}\(I(X;Z) - \beta I(Z;Y)\)

****
\[\min_\rho \left[ I(\rho_\theta; Z) - \beta \cdot I(Z; Y | \rho_\theta) \right]\]

\(I(A;B) = S(\rho_A) + S(\rho_B) - S(\rho_{AB})\)

\textbf{ IB }
\[I(\rho_\theta; Z) \geq S(\rho_\theta) - S(Z|\rho_\theta)\]

\textbf{ D2-1 }****
\[|A|_{quantum} \geq I_{plan}^{quantum} + \delta_{margin}^{quantum}\]


\subsubsection{ 4Schmidt}


\textbf{Schmidt   }
\(|\psi\rangle = \sum_{i=1}^r \lambda_i |i\rangle_A |i\rangle_B\) \(r\)  Schmidt

\textbf{Schmidt -Betti }
\[\beta_1(\mathcal{P}_\theta) \sim r - d\]

\(d\)

\textbf{-}
\[Z[\psi] = \langle\psi|\hat{U}(t)|\psi\rangle = \sum_{C \in \pi_1(\mathcal{M})} \exp(-S[C]/\hbar) \cdot \text{winding}(C)\]


\#\#

\subsubsection{ 1}

```python
\section{1.  PennyLane / Cirq}

\section{2.  S()}

\section{3.    CS() + log(d)}

```

\subsubsection{ 2}

```python
\section{1.}

\section{2.}

\section{3.}

```

\subsubsection{ 3- IB}

```python
\section{1. /}

\section{2.  IB  I(Z;Y)}

```


\#\#

1. \textbf{""}
2. \textbf{ Betti }
3. ****


\#\#

\begin{itemize}
\item [x]  1-4
\item [ ]  1
\item [ ]  2
\item [ ]  3- IB
\end{itemize}



* D5 -*
*2026-05-04 | r01 |  | *

\begin{thebibliography}{99}
\bibitem{ref1} Shannon, C.E. (1948). A Mathematical Theory of Communication.
\bibitem{ref2} Cover, T.M. (2006). Elements of Information Theory.
\end{thebibliography}

\end{document}